%% 第1章 绪论

\chapter{绪论}

\section{选题背景及意义}

\subsection{选题目的}

随着大语言模型（Large Language Model，LLM）推理能力的涌现，人工智能正经历从"被动响应工具"
向"智能代理"（AI Agent）的范式转变。AI Agent 能够主动感知环境、规划路径、调用工具并执行
复杂任务，已成为大模型落地应用的关键技术方向。

为降低 Agent 的开发门槛，业界涌现出多种可视化编排平台。这些平台通过拖拽方式让开发者无需
编写大量胶水代码即可完成 Agent 工作流的设计，典型代表包括 Dify\upcite{dify2024}、Coze\upcite{coze2024}
等低代码平台，以及 LangGraph\upcite{langgraph2024}、AutoGen\upcite{wu2023autogen} 等开发框架。

然而，当前主流的 Agent 编排工具主要基于 Python 或 JavaScript 生态构建，难以与企业现有的
Java 微服务架构深度集成。在大规模、高并发、强类型的企业生产环境中，Java 生态的稳健性不可替代，
但 Java 开发者却缺乏一套原生的、功能完备的 Agent 可视化编排平台。

基于此背景，本研究旨在设计并实现一套\textbf{基于 Spring Boot 的 AI Agent 智能体编排系统}，
该系统融合 Spring AI 框架与可视化 DAG 编排能力，支持声明式配置、人机协作、知识增强与安全访问控制，
为 Java 开发者提供企业级的智能体开发解决方案。

\subsection{选题意义}

\textbf{在理论层面}，本研究探索声明式配置与 DAG 调度相结合的 Agent 编排方法。
当前 Agent 编排系统的研究多集中于 Python 生态，如何在 Java 环境下实现业务流程、知识检索、
人工审核与安全控制的一体化，仍缺乏系统性的工程化探索。本研究通过构建 Java 工程化系统，
为 AI Agent 在企业级场景中的落地提供新的实践参考。

\textbf{在应用层面}，本研究具有以下实用价值：
\begin{enumerate}
  \item 填补 Java 生态空白，使 Java 开发者能够利用现有技术栈进行 AI 应用开发；
  \item 通过可视化编排、DAG 调度和状态持久化机制，提高 AI 系统的可配置性与执行稳定性；
  \item 引入人工检查点、知识检索增强与认证安全机制，提升系统在企业场景下的可控性与可用性。
\end{enumerate}

\section{国内外研究现状}

\subsection{低代码 Agent 平台}

面向非技术用户的低代码平台发展迅速。国外以 Dify\upcite{dify2024}、Coze\upcite{coze2024}
等为代表，提供可视化的 Agent 搭建界面。国内方面，百度推出了 AgentBuilder，
阿里推出了 ModelScope-Agent，腾讯推出了元器等产品，均提供类似的低代码开发能力。

这类平台的优势在于入门门槛低，适合快速原型验证。但其局限性也很明显：
难以与企业现有的后端系统进行深度集成，只能通过 HTTP 端点等网络请求方式进行交互。

\subsection{Python 生态编排框架}

Python 生态处于领先地位。LangChain\upcite{langchain2023} 是目前最流行的 LLM 应用开发框架，
提供丰富的模型接口、工具集成和记忆管理能力，但早期"链式"设计在处理复杂循环逻辑时显得力不从心。

为解决这一问题，LangGraph\upcite{langgraph2024} 引入图论概念，将 Agent 的状态流转建模为节点与边的组合，
天然支持循环、分支和复杂流程控制，成为当前 Agent 工作流编排的重要方案。

微软研究院推出的 AutoGen\upcite{wu2023autogen} 则采用"对话即计算"范式，
让多个 Agent 以对话方式完成任务分解与协作，支持人机协作和代码执行等功能。

上述 Python 框架功能强大，但对 Java 开发者存在较高的学习和集成门槛，
难以直接应用于企业现有的 Spring Boot 微服务架构中。

\subsection{Java 生态现状}

Spring AI\upcite{springai2024} 作为 Spring 官方推出的 AI 应用框架，在近期逐步成熟，
提供了统一的模型接口（ChatClient）、结构化输出和函数调用机制，但目前主要解决 LLM 接入层面的问题，
尚未提供成熟的工作流编排能力。

Java 生态中缺乏一个类似 LangGraph 的、基于配置驱动的图导向编排引擎，
开发者仍需手动编写大量代码来管理 Agent 的状态流转和任务调度。\textbf{这一空白正是本研究的切入点。}

\section{研究目标与内容}

\subsection{研究目标}

本研究的目标是设计并实现一套基于 Spring AI 的 AI Agent 编排系统，
围绕企业级智能体应用的配置、执行、审核、增强与安全支撑需求，形成以下五项核心能力：

\textbf{（1）Agent 管理与可视化编排：}系统支持通过可视化画布拖拽节点、配置参数与连接关系，
以 JSON 图结构保存 Agent 工作流，并支持版本发布、回滚与管理。

\textbf{（2）工作流调度执行：}系统支持基于 DAG 的事件驱动调度，
实现节点并行执行、条件分支、路径剪枝与状态推进。

\textbf{（3）人工审核与恢复机制：}系统支持在关键节点引入人工检查点，
实现暂停、查看上下文、修改中间结果、审批恢复或拒绝终止。

\textbf{（4）知识库与检索增强：}系统支持文档上传、切分、向量化入库与语义检索，
在工作流执行过程中为 LLM 节点注入外部知识上下文。

\textbf{（5）用户认证与系统安全支撑：}系统支持邮箱注册、登录、找回密码和令牌认证，
并通过验证码频率控制、密码加密与访问控制保障系统安全。

\subsection{研究内容}

研究内容分为四个板块：

\textbf{（1）需求分析：}分析 Java 开发者、业务人员、系统管理员等角色的使用场景，
围绕 Agent 管理、工作流执行、人工审核、知识检索与用户认证构建需求模型。

\textbf{（2）系统设计：}采用 DDD 分层架构思想，设计前后端分离与领域驱动相结合的系统结构，
明确可视化编排、执行调度、审核恢复、知识处理与认证安全等模块之间的数据流与职责边界。

\textbf{（3）功能实现：}重点实现 Agent 管理与可视化编排、工作流调度执行、人工检查点与恢复、
知识库与检索增强、用户认证与系统安全支撑五个核心模块。

\textbf{（4）系统测试：}覆盖单元测试、集成测试和功能测试三个层面，
验证五个核心模块的功能闭环与关键链路可用性。

\section{本章小结}

本章阐述了本文的选题背景与研究意义，综述了低代码 Agent 平台、
Python 编排框架和 Java 生态的国内外研究现状，
明确了本文面向 Java 生态的研究定位，并将研究目标统一收束为
Agent 管理与可视化编排、工作流调度执行、人工审核与恢复机制、
知识库与检索增强、用户认证与系统安全支撑五个核心模块，
为后续研究工作奠定了背景基础。
